\chapter*{专题练习3}
填空题：
\begin{enumerate}
    \item 若函数$y = f(x)$的值域是$[\frac{1}{2}, 3]$，则函数$F(x) = f(2x+1) + \frac{1}{f(2x+1)}$的值域是\blkda{$[2, \frac{10}{3}]$}
    \begin{jieda}
        换元即可.
    \end{jieda}
    \item 函数$f(x) = \frac{x^2 + 3x + 2}{x+3} (-2<x<0)$的值域为\blkda[2]{$[2\sqrt{2}-3, \frac{2}{3})$}
    \begin{jieda}
        设$t = x+3 (t \in (1, 3))$，则$y = \frac{t^2-3t+2}{t} = t + \frac{2}{t} - 3$，根据对勾函数性质可知其值域为$[2\sqrt{2}-3, \frac{2}{3})$
    \end{jieda}
    \item 函数$f(x) = \sqrt{\frac{x+1}{x^2+4x+7}}$的值域是\blkda{$[0, \frac{\sqrt{6}}{6}]$}
    \begin{jieda}
        首先换元，令$t = x+1, \sqrt{\frac{x+1}{x^2+4x+7}} = \sqrt{\frac{t}{t^2 + 2t + 4}} =  \left\{ 
         \begin{array}{ll}  
            \sqrt{\frac{1}{t + \frac{4}{t} + 2}}, & t \neq 0, \\ 
            0, & t = 0.
         \end{array}
        \right. $，根据对勾函数的性质可知原函数的值域为$[0, \frac{\sqrt{6}}{6}]$。
    \end{jieda}
    \item 函数$y = \frac{1}{x^2+2x-4}$的严格增区间是\blkda[5.5]{$(-\infty,-1-\sqrt{5})$和$(-1-\sqrt{5}, -1]$}. 
    \begin{jieda}
        \textbf{定义域先行}：$x \in \{x | x \neq -1\pm \sqrt{5}\}$ \\
        外层函数为反比例函数，有两个严格减区间，故要寻找整体的严格增区间，应寻找内层函数的严格减区间即$(-\infty,-1]$ \\
        综合两者，因此是两个单调区间：$(-\infty,-1-\sqrt{5})$和$(-1-\sqrt{5}, -1]$ \\
        \ps 单调区间之间不能用并集符号连接！
    \end{jieda}
    \item 函数$f(x) = (\frac{1}{2})^{\sqrt{x^2-3x-4}}$的递增区间是\blkda[3]{$(-\infty, -1]$}. 
    \begin{jieda}
        \textbf{定义域先行}：$x \in (-\infty, -1] \cup [4, +\infty)$ \\
        外层函数为单调减的指数函数，故要寻找整体的严格增区间，应寻找内层函数的严格减区间即$(-\infty,-1]$
    \end{jieda}
    \item 函数$f(x)=\log_2\left(x-\frac{a}{x}\right)(a>0)$在区间$[2,+\infty)$上是增函数, 则$a$的取值范围是\blkda{$(0,4)$}. 
    \begin{jieda}
        当$a>0$时, $y=x-\frac{a}{x}$在$(0,+\infty)$上单调递增. \\
        考虑到\textbf{定义域}, 于是$2-\frac{a}{2}>0$, 故$a\in(0,4)$. 
    \end{jieda}
    \item 已知函数$y = log_2(x^2 - ax - a)$在区间$(2, +\infty)$上严格增，则$a$的取值范围是\blkda{$(-\infty, \frac{4}{3}]$}. 
    \begin{jieda}
        外层函数严格增，故内层函数$y = x^2 - ax - a$满足在区间$(2, +\infty)$上严格增且函数值大于0.故应有$\frac{a}{2} \leq 2$且$4-3a \geq 0$，故$a \leq \frac{4}{3}$
    \end{jieda}
    \item 已知函数 $f(x)=\frac{\sqrt{2021-a x}}{a-1}$ 在 $[0,1]$ 上严格减，那么实数 $a$ 的取值的范围是\blkda[3]{$(-\infty, 0) \cup(1,2021]$}.
    \begin{jieda}
    首先判断“定义域”，$2021-ax \geq 0$在$x \in [0, 1]$恒成立，设$g(x) = 2021-ax$，这是一条斜率不确定的线段，只需两个端点满足即可恒成立，即$g(0) \geq 0, g(1) \geq 0$，因此$a \leq 2021$，又因为$a \neq 1$，因此$a \in (-\infty, 1) \cup (1, 2021]$时$f(x)$在$[0, 1]$上有定义. \\
    当 $a<0$ 时, $y=\sqrt{2021-a x}$ 在 $[0,1]$ 上严格单调增, 且 $a-1<0$,所以 $f(x)$ 在 $[0,1]$ 上严格单调减, 符合题意, 当 $a=0$ 时, $f(x)=-\sqrt{2021}$ 并非严格增或严格减, 不符合题意, 当 $0<a<1$ 时, $y=\sqrt{2021-a x}$ 在 $[0,1]$ 上严格单调减, 且 $a-1<0$, 不符合题意, 当 $2021\geq a>1$ 时, $y=\sqrt{2021-a x}$ 在 $[0,1]$ 上严格单调减, $a-1>0$, 符合题意.\\
    综上实数 $a$ 的取值的范围是 $(-\infty, 0) \cup(1,2021]$. \\
    \ps \textbf{$y = a \cdot f(x)$，即一个函数乘常系数，类型的函数的单调性在$a > 0$时与$f(x)$保持一致，在$a < 0$时与$f(x)$完全相反.}
    \end{jieda}
    \item 若函数$y=\frac{a^x}{2a^2-5a+2}$在$\bR$上为严格减函数, 则$a$的取值范围是\blkda[3]{$(0,\frac{1}{2})\cup(1,2)$}. 
    \begin{jieda}
    $\begin{cases}
    a\in(0,1)\\
    2a^2-5a+2>0
    \end{cases}$
    或$\begin{cases}
    a>1\\
    2a^2-5a+2<0
    \end{cases}$, 
    解得$a\in(0,\frac{1}{2})\cup(1,2)$. 
    \end{jieda}
    \item 已知函数 $f(x)=\left\{\begin{array}{l}x^{2}+2 a x, x \geq-1 \\ (2-a) x+2 a-5, x<-1\end{array}\right.$, 若 $f(x)$ 在 $\mathbf{R}$ 上是严格增函数, 则实数 $a$ 的取值范围是\blkda{$1 \leq a \leq \frac{8}{5}$} .
    \begin{jieda} 由函数 $f(x)=\left\{\begin{array}{l}x^{2}+2 a x, x \geq-1 \\ (2-a) x+2 a-5, x<-1\end{array}\right.$ 在 $\mathbf{R}$ 上是严格增函数, \\
    得 $\left\{\begin{array}{l}-a \leq-1 \\ 2-a>0 \\ 1-2 a \geq-(2-a)+2 a-5\end{array}\right.$,解得 $1 \leq a \leq \frac{8}{5}$,
    所以实数 $a$ 的取值范围是 $1 \leq a \leq \frac{8}{5}$.
    故答案为: $1 \leq a \leq \frac{8}{5}$ \\
    \ps 分段函数严格单调增等价于各段都严格增，且在端点处相连或者满足“更上一层楼型”
    \end{jieda}
    \item 已知函数$f(x) = \left \{ \begin{array}{ll}
         (2-5a)^x & x < 1  \\
         \frac{a}{x} + 5 & x \geq 1
    \end{array}\right.$ 在$\mathbf{R}$上严格递增，则实数$a$的取值范围是\blkda{$[-\frac{1}{2}, 0)$} .
    \begin{jieda}
        函数$f(x)$在$\mathbf{R}$上严格递增的充要条件为$\left \{ \begin{array}{l}
         2-5a > 1  \\
         a < 0 \\
         a + 5 \geq 2 - 5a
    \end{array}\right.$，解得$a \in [-\frac{1}{2}, 0)$
    \end{jieda}
    \item 对于实数$a, b$定义运算“*”：$a*b = \left\{  
         \begin{array}{ll}  
            a^2 - ab, & a \leq b, \\  
            b^2 - ab, & a > b.
         \end{array}
    \right.$ 设$f(x) = (2x-1) * (x-1)$，且关于$x$的方程$f(x) = m (m \in \mathbf{R})$恰有三个互不相等的实数根$x_1, x_2, x_3$，则$x_1x_2x_3$的取值范围是\blkda{$(-\frac{\sqrt{3} - 1}{16}, 0)$}
    \begin{jieda}
        根据运算定义规则，首先求$2x-1$和$x-1$的交点，为$(0, -1)$，\\
        故$f(x) = \left\{ 
                 \begin{array}{ll}  
                    (2x-1)^2 - (2x-1)(x-1), & x \leq 0, \\  
                    (x-1)^2 - (2x-1)(x-1), & x > 0.
                 \end{array}
                \right. $ \\
        即$f(x) = \left\{ 
                 \begin{array}{ll}  
                    2x^2-x, & x \leq 0, \\  
                    x-x^2, & x > 0.
                 \end{array}
                \right. $ \\
        观察函数图像可知，在满足题意的情况下，$y = m (0 < m < \frac{1}{4})$与$f(x)$在$y$轴左侧有一个交点，在$y$轴右侧有两个交点。故由韦达定理有$x_2x_3 = m$，$x_1 = \frac{1-\sqrt{1+8m}}{4}$。故$x_1x_2x_3 = m \cdot \frac{1-\sqrt{1+8m}}{4} = -\frac{1}{4}(m \cdot (\sqrt{1+8m} - 1))$。因为$-\frac{1}{4}(m \cdot (\sqrt{1+8m} - 1))$在$m \in (0, \frac{1}{4})$单调递减，故其取值范围是$(-\frac{\sqrt{3} - 1}{16}, 0)$
    \end{jieda}
        \item 已知函数$f(x) = \left\{  
         \begin{array}{ll}  
            x^3, & x \geq 0, \\  
            -x, & x < 0.
         \end{array}
    \right.$ 若函数$g(x) = f(x) - |kx^2-2x| (k \in \mathbf{R})$恰有4零点，$k$的取值范围是\blkda[3.5]{$(-\infty, 0) \cup (2\sqrt{2}, +\infty)$}.
    \begin{jieda}
        容易判断，$g(x)$有一零点为0. 设$F(x) = \left\{  
         \begin{array}{ll}  
            x^2, & x \geq 0, \\  
            1, & x < 0.
         \end{array}
    \right.$，$H(x) = |kx-2|$，故其余三个零点是函数$G(x) = F(x) - H(x)$的零点.易知$k = 0$时$G(x)$没有三个零点. 
    \begin{enumerate}
        \item 当$k < 0$时，$G(x)$在$(-\infty, \frac{2}{k})$单调增，在$(\frac{2}{k}, 0)$单调减，在$(0, -\frac{k}{2})$单调减，在$(-\frac{k}{2}, +\infty)$单调增.$G(\frac{3}{k}) = 0, G(\frac{1}{k}) = 0, G(0) = -2, G(2-k) = 2-2k$，故$G(x)$有三个零点，分别是$\frac{3}{k}, \frac{1}{k}$，以及一个位于$(0, 2-k)$的零点.
        \item 当$k > 0$时，当$x < 0$，$G(x) = kx-1$，无零点. \\
        当$x \in (0, \frac{2}{k})$，$G(x)$单调增,$G(\frac{1}{k}) = 0$. \\
        当$x \geq \frac{2}{k}$，$G(x) = x^2 - kx + 2$，要使得$G(x)$在$[\frac{2}{k}, +\infty)$有两个零点，必有$\Delta > 0$，即$k > 2\sqrt{2}$. 抛物线的对称轴为$x = \frac{k}{2}$，因为当$k > 2\sqrt{2}$时$\frac{k}{2} > \frac{2}{k}$，又因为$G(\frac{2}{k}) > 0$，所以当$k > 2\sqrt{2}$时，$G(x)$在$[\frac{2}{k}, +\infty)$有两个零点.
    \end{enumerate}
        综上，$k \in (-\infty, 0) \cup (2\sqrt{2}, +\infty)$ \\
        \ps 从数形结合角度平移$H(x)$来看会更直观.
    \end{jieda}
    \item 已知函数$f(x) = \left\{  
         \begin{array}{ll}  
            x^2-2ax+9, & x \leq 1, \\  
            x+\frac{4}{x}+a, & x > 1.
         \end{array}
    \right.$ 若$f(x)$的最小值为$f(1)$，则实数$a$的取值范围是\blkda{$[2, +\infty)$}
    \begin{jieda}
        首先看到函数在$x \leq 1$时是开口向上的二次函数，函数在$x > 1$时是对勾函数。由于$f(1)$是最小值，所以$f(x)$在$(-\infty, 1]$应该单调递减，故其对称轴应该位于$x=1$及其右侧，即$a \geq 1$。这样我们有$f(1) = 10-2a < f(x), x < 1$。又因为$f(x)$在$(1, +\infty)$为对勾函数，在$(1, 2]$单调递减，在$(2, +\infty)$单调递增，$f(x)$在$(1, +\infty)$的最小值在$x = 2$取得，$f(2) = 4+a$, 故$4+a \geq 10-2a$, 故$a \geq 2$ \\
        \ps 分段函数的最小值要比每一段内所有的函数值都要小.
    \end{jieda}
    \item “$a+b=0$”是“函数$f(x)=x\abs{x+a}+b$是奇函数”的\blkda[2]{必要不充分}条件
        \begin{jieda}
        奇函数若在$x=0$处有定义, 则此处的函数值为零, 即$f(0)=0$, 于是$b=0$. \\
        故$f(x)=x\abs{x+a}$. 于是由$f(x)$为奇函数，不难得到$a=0$, 比如$f(-a)+f(a)=0$. 
        \end{jieda}
        \item 若函数$f(x)=\frac{a-\me^x}{1+a\cdot\me^x}$($a$为常数)在定义域上为奇函数, 则$a$的值为\blkda{$\pm 1$}. 
        \begin{jieda}
        首选解法：$f(0)=0$或无意义，即$\frac{a-1}{1+a} = 0$或$\frac{a-1}{1+a}$无意义，故$a = \pm 1$\\
        也可以采用就极限思想：让函数在$\pm \infty$时的极限互为相反数求解：$x \rightarrow +\infty, f(x) \rightarrow -\frac{1}{a}, x \rightarrow -\infty, f(x) \rightarrow a$，由$a = \frac{1}{a}$得$a = \pm 1$. \\
        把握定义求解有：$f(x) = \frac{a-\me^x}{1+a\cdot\me^x}, f(-x) = \frac{a-\me^{-x}}{1+a\cdot\me^{-x}} = \frac{a\me^{x}-1}{a+\me^{x}}$ \\
        由$f(x) = -f(-x)$知$\frac{a-\me^x}{1+a\cdot\me^x} = \frac{1-a\me^{x}}{a+\me^{x}}$，交叉相乘得到$a^2 - e^{2x} = 1-a^2e^{2x}$，移项整理得$(a^2-1) = (1-a^2)\cdot(e^{2x})$，因为此式恒成立，所以必有$a = \pm 1$
        \end{jieda}
        \item 已知$\mathbf{R}$上的奇函数$f(x)$与偶函数$g(x)$满足$f(x) - g(x) = 2^x$，若对于$\forall x \in (0, 2]$，$mf(x) - g(2x) \geq 0$恒成立，则$m$的取值范围是\blkda[2]{$[-2\sqrt{2}, +\infty)$}
        \begin{jieda}
            $f(x) = \frac{2^x - 2^{-x}}{2}, g(x) = -\frac{2^x+2^{-x}}{2}$ \\
            故$m(\frac{2^x - 2^{-x}}{2}) + \frac{4^x+4^{-x}}{2} \geq 0$在$(0, 2]$上恒成立. \\
            换元$t = 2^x - 2^{-x}, t \in (0, \frac{15}{4}]$，则$4^x+4^{-x} = t^2+2$，故有$mt + t^2 + 2 \geq 0$在$(0, \frac{15}{4}]$上恒成立. \\
            参变分离有$m \geq -(t + \frac{2}{t})$在$(0, \frac{15}{4}]$上恒成立.由对勾函数性质可知$-(t + \frac{2}{t})$在$(0, \frac{15}{4}]$上最大值为$-2\sqrt{2}$，因此有$m \in [-2\sqrt{2}, +\infty)$ \\
            \ps \textbf{一个定义域关于原点对称的函数$f(x)$一定能够拆成一个奇函数和一个偶函数之和：$f(x) = \frac{f(x) - f(-x)}{2} + \frac{f(x)+f(-x)}{2}$}
        \end{jieda}
        \item 已知函数$f(x) = \left \{ \begin{array}{ll}
            x^2 + x & x \leq 0 \\
            ax^2 + bx & x > 0
        \end{array}\right.$为奇函数，则$a+b = $\blkda{0}
        \begin{jieda}
            代入特值有$f(1) = a+b$，$f(-1) = 0$，因为$f(x)$是奇函数，所以有$a+b = 0$ \\
            \ps 根据奇偶性求参数可以考虑代入特值（必要条件），但最后需要检验是否充要（本题只问$a+b$并不涉及）.
        \end{jieda}
\end{enumerate}
\clearpage
解答题：
\begin{enumerate}
    \item 求下列函数的值域: \\
    \begin{enumerate*}
    \item $y=\frac{3+x}{4-x}$; 
    \item $y=\frac{2x^2+3x+8}{x^2+x+4}$; 
    \item $y=\frac{x^4+x^2+5}{(x^2+1)^2}$.
    \end{enumerate*}
    \begin{jieda}
    \begin{enumerate}
    \item $y=-1+\frac{7}{4-x}\in(-\infty,-1)\cup(-1,+\infty)$. 
    
    \item $y=2+\frac{x}{x^2+x+4}$. \\
    当$x = 0$时，$y = 2$；当$x \neq 0$时，$y = 2 + \frac{1}{x + \frac{4}{x} + 1}$.\\
    $x + \frac{4}{x} + 1 \in (-\infty, -3]\cup [5, +\infty)$，故$y \in [\frac{5}{3}, \frac{11}{5}]$
    
    \item $y=\frac{(x^2+1)^2-(x^2+1)+5}{(x^2+1)^2}=1 - \frac{1}{(x^2+1)} + \frac{5}{(x^2+1)^2}$. \\
    令$t = \frac{1}{(x^2+1)}, t \in (0, 1]$，故$y = 5t^2-t+1 \in [\frac{19}{20}, 5]$
    \end{enumerate}
    \end{jieda}
    \begin{vspace_}
    \vsp[6]
    \end{vspace_}
    \item 求下列函数的值域: \\
    \begin{enumerate*}
    \item $y=x+4\sqrt{1-x}$; 
    \item $y=\frac{\sqrt{x^2+1}}{x-1}$. 
    \item $y = \sqrt{1+x} - \sqrt{2+x}$.
    \item $y = \sqrt{1+x} + \sqrt{2-x}$.
    \end{enumerate*}
    \begin{jieda}
    \begin{enumerate}
    \item 令$t=\sqrt{1-x}\in[0,+\infty)$, 则$x=1-t^2$, 故$y=1-t^2+4t=-(t-2)^2+5\in(-\infty,5]$. \\
    \ps 可化为二次的. 
    
    \item 当$x>1$时, $y=\sqrt{\frac{x^2+1}{x^2-2x+1}}=\sqrt{1+\frac{2x}{x^2-2x+1}}=\sqrt{1+\frac{2}{x+\frac{1}{x}-2}}\in(1,+\infty)$. \\
    当$x<1$时, $y=-\sqrt{\frac{x^2+1}{x^2-2x+1}}=-\sqrt{1+\frac{2x}{x^2-2x+1}}$. \\
    再分为$(-\infty,0)\cup\{0\}\cup(0,1)$讨论可得$y\in(-\infty,-\frac{1}{\sqrt{2}}]$. \\
    综上, $y\in(-\infty,-\frac{1}{\sqrt{2}}]\cup(1,+\infty)$.
    \item $y = \sqrt{1+x} - \sqrt{2+x} = \frac{(\sqrt{1+x} + \sqrt{2+x})(\sqrt{1+x} - \sqrt{2+x})}{(\sqrt{1+x} + \sqrt{2+x})} = \frac{-1}{(\sqrt{1+x} + \sqrt{2+x})}$ \\
    易知函数在定义域$[-1, +\infty)$内单调递增，最小值在左端点$x = -1$取$y = -1$，当$x$趋于无穷时，函数值趋于0。因此值域为$[-1, 0)$
    \item $y = \sqrt{1+x} + \sqrt{2-x}$，故$y^2 = (1+x) + (2-x) + 2\sqrt{(1+x)(2-x)}  = 3 + 2\sqrt{(1+x)(2-x)} \in [3, 6]$，故$y \in [\sqrt{3}, \sqrt{6}]$
    \end{enumerate}
    \end{jieda}
    \begin{vspace_}
    \vsp[6]
    \end{vspace_}
    \item 设$a$为实数, 函数$f(x)=2x^2+(x-a)\abs{x-a}$. \\
    \begin{enumerate*}[1.]
    \item 若$f(0)\geq 1$, 求$a$的取值范围; 
    \item 求$f(x)$的最小值. 
    \end{enumerate*}
    \begin{jieda}
    \begin{enumerate}[1.]
    \item $f(0)=-a\abs{a}$，因为$g(x)=-x\abs{x}$在$\bR$上严格单调递, 且$g(-1)=1$. \\
    于是$f(0)=-a\abs{a}\geq 1=g(-1)\iff a\leq -1$.  \\
    \ps 想象一下$g(x)=-x\abs{x}$的图象.
    \item $f(x)=\begin{cases}
    x^2+2ax-a^2, & x\leq a\\
    3x^2-2ax+a^2, & x\geq a
    \end{cases}$. \\
    函数图象由两个二次函数图象的部分拼接而成，且拼接处连续，分析两个开口向上的二次函数的对称轴可知，一个为$x = -a$，一个为$\frac{a}{3}$ \\
    讨论拼接点$x = a$与两个对称轴不同的位置关系下函数的单调性，即可得到函数的最小值，在不同的拼接点条件下寻找全局的最小值即为所求. 
    \begin{dingautolist}{172}
        \item $a = 0$时，三者重合，此时最小值为$0$ 
        \item $a > 0$时，$-a < \frac{a}{3} < a$，故函数$f(x)$在第一段上先减后增，在第二段上单调增，因此最小值在第一段的顶点处取得（$x = -a$），最小值为$-2a^2$
        \item $a < 0$时，$a < \frac{a}{3} < -a$，故函数$f(x)$在第一段上单调递减，在第二段上先减后增，因此最小值在第二段顶点取得（$x = \frac{a}{3}$），最小值为$\frac{2a^2}{3}$
    \end{dingautolist}
    综上，$a < 0$时为$\frac{2a^2}{3}$，$a \geq 0$时，$f(x)$的最小值为$-2a^2$. \\
    \ps 最值需要通过单调性来说明，分段函数分段求，分析各段上函数的单调性.
    另外, 本题的另一个麻烦之处在于函数中的分段点与参数$a$有关, 各段中的``对称轴''(单调区间的端点)亦与参数$a$有关. 
    \end{enumerate}
    \end{jieda}
    %\clearpage
    \item 已知函数$f(x) = x^2 + (x-1)|x-a|$，是否存在实数$a$，使得不等式$f(x) \geq 2x-3$对一切实数$x$恒成立，若存在，求出$a$的范围，若不存在，说明理由.
    \begin{jieda}
        不等式$f(x) \geq 2x-3$对一切实数$x$恒成立即$f(x) - 2x + 3 \geq 0$恒成立. 设$g(x) = f(x) - 2x + 3$，$g(x) = \left \{ \begin{array}{ll}
            (a-1)x+(3-a) & x < a \\
            2x^2 - (a+3)x+(a+3) & x \geq a
        \end{array} \right.$，故应有$g(x) \geq 0$恒成立.
        \begin{dingautolist}{172}
            \item $a > 1$ ，此时$g(x)$在$(-\infty, a)$为单调递增的一次函数，不满足要求.
            \item $a = 1$，此时$g(x) = 2(x < a)$，当$x \geq a$时$g(x)$单调递增，故满足要求.
            \item $a < 1$，此时$g(x)$在$(-\infty, a)$为单调递减的一次函数，在$[a, \frac{a+3}{4})$单调减，在$[\frac{a+3}{4}, +\infty)$单调增，故$g(x) \geq g(\frac{a+3}{4}) = (a+3) - \frac{(a+3)^2}{8}$，当$(a+3) - \frac{(a+3)^2}{8} \geq 0$时满足题意，解得$a \in [-3, 1)$
        \end{dingautolist}
        综上，$a \in [-3, 1]$ \\
        \ps 函数解析式内有部分绝对值的情况通常要写成分段函数来分析
    \end{jieda}
    
    \begin{vspace_}
    \vsp[6]
    \end{vspace_}
    \item 函数$f(x) = \frac{x^2+(3a+1)x+c}{x+a} (a, c \in \mathbf{R})$
        \begin{enumerate}
            \item 当$a = 0$时，是否存在实数$c$使得$f(x)$是奇函数
            \item 函数$f(x)$的图像过点$(1, 3)$，且$f(x)$的图像和$x$轴负半轴有两个交点，求实数$a$的取值范围
        \end{enumerate}
        \begin{jieda}
            \begin{enumerate}
            \item 当$a = 0$，$f(x) = \frac{x^2+x+c}{x} = x + 1 + \frac{c}{x} (x \neq 0)$， $f(1) + f(-1) = 2 \neq 0$，故不存在$c$使得$f(x)$为奇函数.
            \item $f(1) = \frac{2+3a+c}{1+a} = 3$, 故$c = 1, f(x) = \frac{x^2+(3a+1)x+1}{x+a}$. 函数的零点是函数的分子部分的零点子集，故$g(x) = x^2+(3a+1)x+1 (x \neq -a)$与$x$轴负半轴有两个交点$x_1, x_2$。
            \begin{equation*}
            \left\{ 
                \begin{array}{l}  
                    \Delta = (3a+1)^2 - 4 > 0, \\ 
                    x_1 + x_2 = -(3a + 1) < 0, \\
                     a^2 - (3a+1)a + 1 \neq 0.
                 \end{array}
            \right \quad
            \end{equation*}
            解得$a > \frac{1}{3} \mbox{且} a \neq \frac{1}{2}$
        \end{enumerate}
        \ps 否定一个函数是奇函数可以通过举反例.
        \end{jieda}
        
    \begin{vspace_}
    \vsp[6]
    \end{vspace_}

    \item 设函数$f(x) = log_a(x-2a) + log_a(x-3a)$，其中$a > 0 \mbox{且} a \neq 1$。若在区间$[a+3, a+4]$上$f(x) \leq 1$恒成立，求$a$的取值范围.
    \begin{jieda}
        由定义域的要求可知$a+3 > 3a$，即$a < \frac{3}{2}$。$f(x)$的单调性由$a$决定，下面进行分类讨论。
        \begin{enumerate}
            \item $\frac{3}{2} > a > 1$ \\
            此时$f(x)$单调递增，故$f(a+4) \leq 1$，即$log_a((4-a)(4-2a)) \leq 1$，\\
            即$(4-a)(4-2a) \leq a$，解得$a \in (\frac{13-\sqrt{41}}{4}, \frac{13+\sqrt{41}}{4})$(舍弃)
            \item $1 > a > 0$ \\
            此时$f(x)$单调递减，故$f(a+3) \leq 1$，即$log_a((3-a)(3-2a)) \leq 1$，\\
            即$(3-a)(3-2a) \geq a$，解得$a \in (-\infty, \frac{5-\sqrt{7}}{2}) \cup (\frac{5+\sqrt{7}}{2}, +\infty)$，结合$1 > a > 0$，得到$a \in (0, 1)$
        \end{enumerate}
        综上，$a \in (0, 1)$\\
        \ps 也可以在分类讨论的时候不解一元二次不等式
        \begin{enumerate}
            \item $\frac{3}{2} > a > 1$ \\
            由于$y = (4-a)(4-2a) - a = 2a^2-13a+16$在$(1, \frac{3}{2})$内单调递减，而$a = \frac{3}{2}$时不等式不成立，因此这种情况无解。
            \item $1 > a > 0$ \\
            由于$y = (3-a)(3-2a) - a = 2a^2-10a+9$在$(0, 1)$内单调递减，而$a = 1$时不等式成立，因此这种情况解集为$(0, 1)$。
        \end{enumerate}
    \end{jieda}
\end{enumerate}